\section{Common Crawl}

Common Crawl is the largest freely available public archive of web crawl data, totalling more than 10 petabytes across crawls published approximately monthly, each typically containing more than two billion web pages \citep{baackCriticalAnalysisLargest2024,CommonCrawlGet}. Widely used as a web-scale source of language data for corpus construction, NLP research, and large language model pretraining \citep{baackCriticalAnalysisLargest2024,wenzekCCNetExtractingHigh2019}, it has been cited in over 12,000 research papers \citep{CommonCrawlGet}. For the present study, its central value lies in breadth: rather than drawing on a single platform, newspaper archive, or curated corpus, Common Crawl provides repeated cross-sections of general web discourse. The archive stores raw HTML and, importantly, collects only a sample of pages from each domain it visits---meaning, for instance, that only a subset of Wikipedia articles will appear, not the full site. Domain selection and page depth are governed by \textit{harmonic centrality}, a graph-based scoring method adopted in 2017 whereby a domain's importance is determined by the volume of direct and indirect inbound links from other domains, with direct links weighted most heavily; higher-scoring domains are both more likely to be crawled and to have more of their pages fetched \citep{baackCriticalAnalysisLargest2024}. Despite its enormous size, Common Crawl is neither a complete copy of the web nor a representative sample of it. A growing number of rights holders---including major outlets such as the New York Times---now block Common Crawl via the \texttt{robots.txt} standard, largely in response to AI training data concerns, and large social media platforms such as Facebook have done so for considerably longer \citep{baackCriticalAnalysisLargest2024}. The data used in this study should therefore not be treated as a representative survey of the web or of public opinion; they reflect what Common Crawl crawled, retained, and made available within these structural constraints.

\Needspace{0.55\textheight}
\subsubsection*{Common Crawl Collection Pipeline}

\begin{figure}[!htbp]
  \centering
  % Source: ../reports/figures/commoncrawl_pipeline/commoncrawl_collection_pipeline.tex
  \includegraphics[width=\textwidth]{commoncrawl_pipeline/commoncrawl_collection_pipeline.pdf}
  \caption{Common Crawl collection pipeline used to construct the trend and corpus materials.}
  \label{fig:commoncrawl-pipeline}
\end{figure}
\FloatBarrier

To collect the data for this study, a Common Crawl collection pipeline was developed to extract quality-gated general discourse for the analysis of lexical semantic change (LSC) in specific target terms, here ADHD and autism, against matched baseline terms. The pipeline workflow is illustrated in Figure~\ref{fig:commoncrawl-pipeline}. Target and baseline terms are processed together so that yearly denominators, sampling logic, and quality filters remain comparable across term groups. The pipeline uses Common Crawl's two main file formats sequentially. WET files provide compact extracted plaintext and are therefore used for large-scale term scanning and yearly prevalence denominators, but they lack the HTML structure and metadata needed for stronger boilerplate filtering. Candidate documents are therefore resolved to their corresponding WARC records, which preserve the full HTML code. Although WARC processing is slower, it enables main-text extraction, boilerplate removal, and metadata recovery. This WET-first, WARC-second design keeps the pipeline economical by reserving expensive WARC processing for candidate documents only.\footnote{For a detailed record of the design, configuration, and methodological decisions behind the Common Crawl collection pipeline refer to \url{https://github.com/jako6f/msc-nlp-therapy-speak/blob/main/reports/commoncrawl_corpus_design_and_provenance.md}.}

The design consists of two linked tracks. The trend track uses fixed-effort annual samples to estimate how frequently target and baseline terms appear over time. The corpus track builds a larger, quality-gated document corpus for downstream NLP analysis. As an additional safeguard against overrepresentation by large websites beyond Common Crawl’s restrictive sampling approach, domain caps are applied at 50 WET-validated hit rows per registered domain per Common Crawl crawl. Intermediate summaries and manifests are retained so that each year, crawl, track, and batch remains auditable. The pipeline is designed to run end-to-end on AWS EC2, using S3 as the durable storage and transfer layer for intermediate and final collection artefacts. Yearly crawl selection is deterministic: one Common Crawl snapshot is selected per year near a fixed annual anchor date and then frozen in a crawl map.

Several software choices are methodologically consequential because they affect corpus membership. WET and WARC records are parsed with \texttt{warcio}, WARC pointers are resolved through a local \texttt{pywb}-based index server, archived HTML is converted to main text with Trafilatura and Resiliparse, and post-extraction filtering uses DataTrove quality filters followed by English-language filtering with \texttt{py3langid} \citep{bevendorffElasticChatNoirSearch2018,barbaresiTrafilaturaWebScraping2021,penedoDataTroveLargeScale2026,luiLangidpyOfftheshelfLanguage2012}.\footnote{Both \texttt{warcio} and \texttt{pywb} are open-source Webrecorder projects; see \url{https://github.com/webrecorder/warcio} and \url{https://github.com/webrecorder/pywb}.}

The data collection spans 13 annual Common Crawl snapshots from 2014 to 2026. This window maximises temporal depth while remaining compatible with a stable WET-first, WARC-second workflow. By 2014, Common Crawl provided sufficiently large WET/WARC-format crawls for efficient plaintext scanning, denominator construction, and HTML-based validation; earlier data are less comparable and require additional handling due to older archive formats. The 2026 endpoint reflects the latest collection year available for the project. In the trend track, the pipeline scanned 55.6 million web pages and retained 156,189 pages featuring validated term hits. In the corpus track, it scanned 220 million web pages and retained 336,178 analysis-ready documents. Of these, 87,173 documents contain ADHD or autism target terms. Target membership is non-exclusive: 31,354 documents contain ADHD terms, 67,614 contain autism terms, and 11,795 contain both. The collection was run on an AWS \texttt{m7i-flex.large} instance, featuring 2 vCPUs, 8 GiB of RAM, and up to 12.5 Gbps network bandwidth \citep{CloudComputingServices}. On this instance type, corpus throughput was approximately one hour per million scanned WET records.

\section{Target Terms}

Two target concepts were selected for diachronic semantic analysis: \textit{ADHD} and \textit{autism}. Target documents were retrieved using the matching expressions shown in Table~\ref{tab:target-patterns}. The abbreviation \texttt{ASD} (autism spectrum disorder) was retained only when \textit{autism} occurred within \(\pm 200\) characters, reducing false positives from unrelated acronym use. The acronym \texttt{ADD} (attention deficit disorder) was not included as a matching expression because ADHD is the current canonical diagnostic label, whereas ADD is older and colloquial \citep{americanpsychiatricassociationDiagnosticStatisticalManual2013}; more importantly, \texttt{ADD} would create substantial false positives in WET scanning because it overlaps with the common verb \textit{add} and with web chrome such as ``add to cart''. The broader expression \texttt{attention[-\string\s]?deficit} retains coverage of relevant expanded forms while avoiding this high-noise acronym.

\begin{table}[htbp]
  \centering
  \small
  \caption{Target-term matching expressions.}
  \label{tab:target-patterns}
  \begin{tabular}{@{}ll@{}}
    \toprule
    Concept & Matching expressions \\
    \midrule
    ADHD & \texttt{\string\badhd\string\b}; \texttt{attention[-\string\s]?deficit} \\
    Autism & \texttt{\string\bautism\string\b}; \texttt{\string\bautistic\string\b}; \texttt{autism[-\string\s]?spectrum}; \texttt{\string\bASD\string\b} \\
    \bottomrule
  \end{tabular}
\end{table}

For comparison, three negative, non-clinical emotion terms were selected with sufficient coverage and interpretable usage: \textit{frustration}, \textit{sadness}, and \textit{loneliness}. These terms are not exact semantic controls for ADHD and autism; they provide a baseline for separating target-specific change from broader shifts in negative affective language in the corpus.

\section{Preprocessing}

Following \citet{baesMultidimensionalFrameworkEvaluating2024}, preprocessing was organised around analysis-specific corpus representations rather than a single all-purpose text file. The downstream semantic analyses use a shared mention-level context table built from the extraction pipeline's corpus product. Target and baseline mentions were re-detected using the frozen collection patterns (see Table~\ref{tab:target-patterns}), including the ASD disambiguation rule, and overlapping matches within the same analysis unit were resolved by keeping the longest span. This prevents expressions such as \textit{autism spectrum} from being counted twice as both a phrase and a shorter nested form. Documents containing both ADHD and autism terms were allowed to contribute to both target groups, with separate mention-level contexts emitted for each concept. Same-sentence acronym-and-expansion pairs occurring within a short local window were also collapsed, so cases such as ``attention deficit hyperactivity disorder (ADHD)'' contribute one conceptual ADHD context rather than two near-duplicate rows. To limit domination by repetitive documents while preserving repeated-use signal, each document could contribute at most three mentions per analysis unit. The semantic time variable was then defined as publication year and only contexts with parseable publication dates between 2014 and 2026 were retained for the shared LSC table. This yielded 293,670 mention contexts from 212,651 documents and 135,771 registered domains, including 28,611 ADHD contexts, 68,253 autism contexts, and separate baseline series for frustration, sadness, and loneliness. For each retained mention, the table stores the matched form, raw-form collapse diagnostics, registered domain, publication and crawl metadata, a \(\pm 5\)-token window for affective collocate analyses, and target-sentence passages for breadth and thematic analyses.

\section{NRC--VAD Lexicon}

The affective analyses use the NRC Valence, Arousal, and Dominance (VAD) Lexicon v2.1 \citep{mohammadNRCVADLexicon2025}, rather than the Warriner norms used in SIBling \citep{baesMultidimensionalFrameworkEvaluating2024}. NRC--VAD offers several advantages for the present study. It is substantially larger, providing human ratings for approximately 45,000 English words and 10,000 multi-word phrases, compared with 13,915 words and no phrase entries in Warriner norms. It has also been reported to have higher reliability, with an aggregate reliability estimate of 0.923 across the three VAD dimensions, compared with 0.823 for the Warriner norms.

NRC--VAD scores range from \(-1\) to \(1\). Higher valence scores indicate more positive affect, higher arousal scores indicate greater emotional activation, and higher dominance scores indicate greater perceived control or power. This study uses valence to estimate whether the local contexts of target terms become more positive or negative over time, and arousal to estimate changes in emotional severity. Dominance is included in the source lexicon but is not used in the present analyses. The ratings were produced using Best--Worst Scaling, in which annotators are shown four items and asked to identify the item that best and worst represents the property being rated \citep{mohammadNRCVADLexicon2025}.
